\section{Roofline of Representative Settings}
\label{sec:predicted}

\paraq{What does this section do?}
It draws the roofline of Section~\ref{sec:roofline} for a few settings taken from the exhaustive
sweep of Section~\ref{sec:hybrid-empirical}, places every layer of every strategy on it, and
looks for the patterns that a configuration algorithm could later exploit. The section is
exploratory: the plots and their presentation are still being tried out, and the algorithm is
left for later (Section~\ref{sec:pred-algo}).

% ---------------------------------------------------------------------------
\subsection{Choice of settings}
\label{sec:pred-settings}

\paraq{Which settings are plotted, and why these?}
We pick three settings, one per winning regime of the measured sweep, so that the plots show
different bottlenecks: one where \textbf{gsplit} is fastest, one where a \textbf{data-parallel
hybrid} with an interior split point $0<k<L$ is fastest, and one where \textbf{data parallelism}
(DGL) is fastest. Pipeline rows marked $^{\dagger}$ in Section~\ref{sec:hybrid-empirical} are not
used for the selection, because of the micro-batch caveat.
% Candidate settings (to confirm with author), chosen so that each pair differs in one knob;
% values are min duration (s) per epoch from results/sweep_all.csv:
%   S1 products, GAT,  L=3, h=1024, NVLink   : DGL 15.48, gsplit 13.09, DP-hybrid k=2 13.63 -> gsplit
%   S2 products, GAT,  L=3, h=1024, PCIe-emu : DGL 15.25, gsplit 16.48, DP-hybrid k=2 14.24 -> hybrid
%   S3 products, SAGE, L=3, h=1024, PCIe-emu : DGL  4.76, gsplit  8.24, DP-hybrid k=2  5.54 -> DGL
%   S1 vs S2 change only the interconnect; S2 vs S3 change only the model.

\begin{table}[h]
\centering
\caption{Settings plotted in this section, with the measured per-epoch time (s) of each
strategy from Section~\ref{sec:hybrid-empirical}. Default 4$\times$A100, fanout 15, batch 1024.
\placeholder{} = to fill.}
\label{tab:pred-settings}
\small
\begin{tabular}{l l l c c l r r r l}
\toprule
ID & Graph & Model & $L$ & $h$ & Interconnect & DGL & gsplit & best hybrid ($k$) & Winner \\
\midrule
S1 & \placeholder & \placeholder & \placeholder & \placeholder & \placeholder
   & \placeholder & \placeholder & \placeholder & gsplit \\
S2 & \placeholder & \placeholder & \placeholder & \placeholder & \placeholder
   & \placeholder & \placeholder & \placeholder & hybrid \\
S3 & \placeholder & \placeholder & \placeholder & \placeholder & \placeholder
   & \placeholder & \placeholder & \placeholder & DGL \\
\bottomrule
\end{tabular}
\end{table}

% ---------------------------------------------------------------------------
\subsection{Inputs}
\label{sec:pred-inputs}

\paraq{What has to be measured to place the points?}
For each setting and each of the three strategies---data parallel ($k{=}L$), gsplit ($k{=}0$),
and the best data-parallel hybrid ($0<k<L$)---one \emph{sampling-only} pass (no features, no
training) records, per layer $\ell$ and GPU $q$, the counts $n_{\ell,q}$, $e_{\ell,q}$,
$m_{\ell,q}$ and the distinct rows $u_\ell$ of Section~\ref{sec:cm-point}. The pipeline needs no
separate pass: its point is computed from the gsplit counts by Eq.~\eqref{eq:cm-pipe-point}.

\paraq{Which rates are used?}
The model constants of Table~\ref{tab:cm-constants} and the measured rates of
Table~\ref{tab:cm-microbench}. Until the microbenchmarks exist, provisional rates are used and
marked as such in each figure: \placeholder.

\begin{table}[h]
\centering
\caption{Per-layer counts on the busiest GPU, from sampling-only runs, and the resulting point
(Eq.~\eqref{eq:cm-point}). One block per setting; shown for S1. \placeholder{} = to fill.}
\label{tab:pred-counts}
\small
\begin{tabular}{l c r r r r r r r r}
\toprule
Strategy & $\ell$ & $n$ & $e$ & $m$ & $u$ & $\delta$ & $r$ & $x$ & $y$ \\
\midrule
data parallel & 0 & \placeholder & \placeholder & 0 & \placeholder & \placeholder & 0 & \placeholder & \placeholder \\
              & 1 & \placeholder & \placeholder & 0 & \placeholder & \placeholder & 0 & \placeholder & \placeholder \\
              & 2 & \placeholder & \placeholder & 0 & \placeholder & \placeholder & 0 & \placeholder & \placeholder \\
\midrule
gsplit        & 0 & \placeholder & \placeholder & \placeholder & \placeholder & \placeholder & \placeholder & \placeholder & \placeholder \\
              & 1 & \placeholder & \placeholder & \placeholder & \placeholder & \placeholder & \placeholder & \placeholder & \placeholder \\
              & 2 & \placeholder & \placeholder & \placeholder & \placeholder & \placeholder & \placeholder & \placeholder & \placeholder \\
\midrule
best hybrid ($k{=}$\placeholder) & 0 & \placeholder & \placeholder & \placeholder & \placeholder & \placeholder & \placeholder & \placeholder & \placeholder \\
              & 1 & \placeholder & \placeholder & \placeholder & \placeholder & \placeholder & \placeholder & \placeholder & \placeholder \\
              & 2 & \placeholder & \placeholder & \placeholder & \placeholder & \placeholder & \placeholder & \placeholder & \placeholder \\
\bottomrule
\end{tabular}
\end{table}

% ---------------------------------------------------------------------------
\subsection{Plots}
\label{sec:pred-plots}

\paraq{What does each plot show?}
One roofline per setting, on the axes of Section~\ref{sec:cm-axes}:
\begin{itemize}\itemsep1pt
  \item roofs: dense and gather (dashed), the binding compute roof (solid), and
        $\hat{B}_{\mathrm{net}}$ for the setting's interconnect (flat);
  \item the regions of Section~\ref{sec:cm-regions}, shaded, with the ridge $x^*(r)$ marked;
  \item one marker shape per strategy (data parallel, gsplit, best hybrid), labelled by layer
        $\ell$;
  \item the pipeline position above each gsplit point (Eq.~\eqref{eq:cm-pipe-point}), drawn as an
        open marker with a vertical line showing the gain of Eq.~\eqref{eq:cm-pipe-gain}.
\end{itemize}

\begin{figure}[h]
  \centering
  \fbox{\parbox{0.85\linewidth}{\centering\small
  \expdiag{roofline of S1 with per-layer points of data parallel, gsplit, and the best hybrid.}\\[4pt]
  \placeholder}}
  \caption{Roofline of setting S1 (\placeholder).}
  \label{fig:pred-s1}
\end{figure}

\begin{figure}[h]
  \centering
  \fbox{\parbox{0.85\linewidth}{\centering\small
  \expdiag{roofline of S2, same layout as Figure~\ref{fig:pred-s1}.}\\[4pt]
  \placeholder}}
  \caption{Roofline of setting S2 (\placeholder).}
  \label{fig:pred-s2}
\end{figure}

\begin{figure}[h]
  \centering
  \fbox{\parbox{0.85\linewidth}{\centering\small
  \expdiag{roofline of S3, same layout as Figure~\ref{fig:pred-s1}.}\\[4pt]
  \placeholder}}
  \caption{Roofline of setting S3 (\placeholder).}
  \label{fig:pred-s3}
\end{figure}

\paraq{Which presentations are being tried?}
\begin{enumerate}\itemsep1pt
  \item \textbf{Single panel per setting}: all strategies and layers together (Figures
        above).
  \item \textbf{Arrows per layer}: an arrow from the gsplit point to the data-parallel point of
        the same layer, so the direction (up or down) shows which strategy wins the layer.
  \item \textbf{Small multiples per layer}: one panel per $\ell$, so the change of region with
        depth is visible.
\end{enumerate}
\placeholder{} (choice of presentation).

% ---------------------------------------------------------------------------
\subsection{Observations}
\label{sec:pred-obs}

\paraq{What do the plots show?}
\placeholder{} (to be written once the plots exist: which region each layer falls in, per
strategy and setting; whether the region agrees with the measured winner of
Table~\ref{tab:pred-settings}).

\paraq{Open question: does the roofline explain the hidden-size trend?}
In Section~\ref{sec:hybrid-empirical} the hybrid's gain over gsplit grows with $h$ in $7$ of $8$
hidden-size ladders. In the roofline, larger $h$ moves points right ($x \propto h$), towards the
compute-bound region where gsplit is favoured (Section~\ref{sec:cm-regions}); with a fixed
$\hat{P}_{\mathrm{bind}}$ this predicts the opposite trend. Whether a $\hat{P}$ that rises with
$h$ (Table~\ref{tab:cm-microbench}), or a change in $\delta$ or $r$, accounts for it is to be
checked on the plots. \placeholder

% ---------------------------------------------------------------------------
\subsection{Towards an algorithm}
\label{sec:pred-algo}

\paraq{What algorithm do the plots suggest?}
\placeholder{} (to be designed from the observations: a rule that maps each layer's region to a
strategy, a split point $k$, and whether to enable the pipeline).

\paraq{Conclusion of this section.}
\placeholder
